\documentclass[11pt,letterpaper]{article}
\usepackage{nsfproposal}
\usepackage{bm}
\usepackage{amsthm}
\usepackage{color,xcolor,colortbl} 
\usepackage{hyperref} 
\usepackage{graphicx} 
\usepackage{vmargin} 
\usepackage{enumitem, soul}
\usepackage{amsmath,amssymb,xspace, comment}
\usepackage{tikz}
\usepackage{multirow}
\usepackage{graphicx}
\usepackage{booktabs,wrapfig}
\usepackage{makecell}
\usepackage{array}
\usepackage{longtable}
\usepackage{tabularray}
\usepackage{tcolorbox}
\usepackage{mfirstuc}
\usepackage{makecell}
\usepackage{algorithm}
\usepackage{algpseudocode}
\usepackage{rotating}
\usepackage{tabularray}
\usepackage{wasysym}

\MFUnocap{are}
\MFUnocap{or}
\MFUnocap{etc}
\MFUnocap{and}
\MFUnocap{with}
\MFUnocap{of}
\MFUnocap{a}


\newtcolorbox{goalsbox}[2][]
{
  colframe = #2!25,
  colback  = #2!10,
  coltitle = #2!20!black,  
  #1,
}

\setpapersize{USletter}
\setmarginsrb{1in}{1in}{1in}{1in}{0pt}{0mm}{0pt}{0mm}


\renewcommand*{\thefootnote}{\fnsymbol{footnote}}
\newcommand{\pseudodot}{{\lower 2.4pt\hbox{$\cdot$}}}
\renewcommand{\title}[1]{\begin{center}\LARGE\bfseries{#1}\end{center}}
\setul{0.8pt}{}
\sodef\an{}{-.02em}{0.2em plus0.2em}{0.5em plus.1em minus.1em}

\newcommand*\cib[1]{\tikz[baseline=(char.base)]{
                \node[shape=circle,fill=black,text=white,draw,inner sep=0.3pt] (char) {#1};}}

\hypersetup{colorlinks=true,citecolor=blue,linkcolor=blue,urlcolor=blue}
\newcommand{\etal}{{\em et al.}\xspace}
\newcommand{\eg}{{\em e.g.,}\xspace}
\newcommand{\ie}{{\em i.e.,}\xspace}
\newcommand{\etc}{{\em etc.}\xspace}
% \renewcommand{\sectionautorefname}{{\S\hspace{-0.5ex}}}
% \renewcommand{\subsectionautorefname}{{\S\hspace{-0.5ex}}}
% \renewcommand{\subsectionautorefname}{Table}
\newcommand{\algorithmautorefname}{Algorithm}
\renewcommand{\tableautorefname}{Table}
\renewcommand{\itemautorefname}{Table}
\renewcommand{\equationautorefname}{Eq.\hspace{-0.5ex}}
\newcommand{\argmin}{\mathop{\mathrm{arg\,min}}}
\newcommand{\argmax}{\mathop{\mathrm{arg\,max}}}

\newcommand{\xo}{\raisebox{0.4ex}{\textbf{$\chi$}}\xspace}
\newcommand{\xod}{{{\xo}-embodied}\xspace}
\newcommand{\xot}{{{\xo}-embodiment}\xspace}
\newcommand{\ours}{{\small$\textsf{{X}-RSecure}$}\xspace}
\newcommand{\BfPara}[1]{{\noindent\bfseries\an{#1}.}\xspace}
\newcommand{\BfParaUL}[1]{{\noindent\ul{\textbf{#1}.}}\xspace}
\newcommand{\dd}{\Circle}
\newcommand{\vv}{\CIRCLE}
\newcommand{\bb}{\RIGHTcircle}
\newcommand{\ee}{$\cdot$}
\newcommand{\BfParaC}[1]{{\vspace{1ex}\noindent\bfseries{\capitalisewords{#1}.}\xspace}}

\newcommand{\todo}[1]{\textcolor{blue}{[$\rightarrow$#1]}}
\newcommand{\mo}[1]{\textcolor{red!85!black}{#1}}
\usepackage{txfonts}
\renewcommand{\rmdefault}{txr}
% \usepackage{algpseudocode}
% \usepackage[linesnumbered]{algorithm2e}

\usepackage{tcolorbox}
\tcbuselibrary{theorems}


% 


\newtcbtheorem[]{observation}{\textbf{}}%
{colframe=white!95!black,fonttitle=\bfseries,coltitle=black, boxsep = 2pt, left = 0pt, right = 0pt, top = 1pt, bottom = 2pt, boxrule = 0pt, bottomrule = 0pt, toprule = 0pt}{th}{}


\newtheoremstyle{styledef}%
{\topsep}% Space above
{\topsep}% Space below 
{}% Body font
{3pt}% Indent amount
{\bfseries}% Theorem head font
{.}% Punctuation after theorem head
{3pt}% Space after theorem head
{}% Theorem head spec (can be left empty, meaning ‘normal’)
\theoremstyle{styledef}

%\newtheorem{observation}{Observation}

\newcommand{\observationdef}[2][0.97\textwidth]{
  \vspace{2ex}\par\noindent\tikzstyle{mybox} = [fill=gray!10,
   thick,rectangle,inner sep=6pt,path picture={\fill [white!40!black] ([xshift=-8.1cm]path picture bounding box.north) rectangle (path picture bounding box.south west);}]
  \begin{tikzpicture}
  \centering
   \node [mybox] (box){%
    \begin{minipage}{#1}{ #2}\end{minipage}
   };
  \end{tikzpicture}
}

\setlength{\baselineskip}{12.6pt} % in text mode
\setlength{\normalbaselineskip}{12.6pt} % in math mode

\usetikzlibrary{positioning}


% Make all figure and table captions small
\usepackage{caption}
\captionsetup[figure]{font=small}
\captionsetup[table]{font=small}




\usepackage{bm}
\usepackage{color,xcolor,colortbl} 
\usepackage{hyperref} 
\usepackage{graphicx} 
\usepackage{vmargin} 
\usepackage{enumitem, soul}
\usepackage{amsmath,amssymb,xspace, comment}
\usepackage{tikz}
\usepackage{multirow}
\usepackage{graphicx}
\usepackage{booktabs,wrapfig}
\usepackage{makecell}
\usepackage{array}
\usepackage{longtable}
\usepackage{tabularray}
\usepackage{tcolorbox}

\begin{document}
%
% Remove page numbers---Research.gov will paginate the document instead
%
\pagestyle{empty} 
%
% Line spacing: 
% The Proposal and Award Policies and Procedures Guide
% (PAPPG: https://www.nsf.gov/publications/pub_summ.jsp?ods_key=pappg)
% mandates, in Chapter 2, section B.2, that the main text should have no more than
% six lines to the vertical inch. With 72 points per inch, the minimum line skip 
% would be 12 points.
%
\setlength{\baselineskip}{12.6pt} % in text mode
\setlength{\normalbaselineskip}{12.6pt} % in math mode



\title{\Large
%Robust Human-in-the-Loop Robotic Learning Using Large Vision-Language Models
SaTC 2.0: RES: Toward Secure and Robust Generalist Robotic Models}
\begin{center}
   {\bfseries\filcenter
Mohammed Abuhamad and George K. Thiruvathukal,  Loyola University Chicago} 
\end{center}



\section{Overview}

This proposal aims to address the security and safety challenges in developing foundation models for robotics, ensuring robust frameworks and defense mechanisms for such models as they incorporate various pre-trained models for different modalities, and securing against inherited vulnerabilities and adversarial threats.
Generalist foundation robotic models are inspired by advancements in large language, vision, and vision-language models and represent a pivotal development in intelligent robotic systems. 
Cross-embodiment robotic datasets and new end-to-end training methods are becoming more democratized, paving the way for significant advancements in foundation robotic models in the near future.
% These models incorporate various pre-trained models (\eg vision and language models) that are trained on Internet-scale datasets.
% Recent studies show that these models enable robotic systems to reason, interact, and learn new tasks in diverse unseen environments.
This proposal focuses on creating comprehensive frameworks, benchmarks, advanced security measures, and specialized methods to enhance the security and reliability of these models for practical applications.
% This main objective can be achieved through three main tasks: \cib{1} performing a comprehensive vulnerability analysis in foundation robotic models, \cib{2} investigating the modality-inherited vulnerabilities and their impact on multimodal robots, and \cib{3} developing benchmarks to ensure the robustness and trustworthiness of robotic models.
% To successfully conduct these tasks, we must overcome unique challenges of building and examining robotic models due to their complex architectures, the multimodal nature of their input, restrictions on the output/action space (which depends on the robot embodiment and data), and the evolving security landscape.

\BfPara{Keywords} trustworthy AI/ML security; robotic foundation models; vision-language-action models; adversarial robustness; multimodal security; safety benchmarking.


\section{Intellectual Merit}
Robotic foundation models carry an attack surface absent from vision and language models. Action heads quantize continuous control into 256 uniform bins, policies emit chunks of consecutive actions, and every predicted token maps to a physical displacement bounded by joint limits and the reachable workspace. Prior security work adapts pixel-space attacks from VLMs and scores success in token space, leaving physical consequence unmeasured. This project grounds attacks, defenses, and measurement in the action space of the robot, delivering \cib{1} a threat taxonomy for multimodal robotic systems across embodiments and tasks, \cib{2} methods for analyzing cross-modal attack propagation, and \cib{3} standardized benchmarks for robotic security, safety, and interpretability. Five artifacts realize these contributions.
\cibx{1} Three attack families exploit structure specific to VLAs. Action-Bin Boundary Perturbation (ABBP) drives action tokens across quantization boundaries under Jacobian-weighted feasibility constraints. Surrogate-Transfer Action Perturbation (STAP) transfers perturbations from surrogates under controlled backbone, data, and embodiment mismatch. The Temporal Trajectory Drift Attack (TTDA) accumulates trajectory drift under a query budget while every step stays kinematically plausible. Poisoning and privacy variants extend these families across the NIST Trustworthy and Responsible AI attack taxonomy.
\cibx{2} Two defenses operate on predicted actions rather than inputs, the Temporal Action Consistency Detector (TACD) and Action-Grounded Adversarial Fine-Tuning (AGAT), so both hold when the compromised modality is unknown. 
\cibx{3} One evaluation protocol reports task success under attack, end-effector deviation, and the perturbation budget needed to halve task success, replacing the sole reliance on action-token-deviation.
\cibx{4} Two interpretability methods explain robot behavior under attack, modality attribution isolating each modality contribution to an adversarial objective by selective reversion, and Spatial-Action Attribution (SAA) tracing a predicted action to image regions and instruction tokens.
\cibx{5} VLA-SecBench releases the model zoo, hardened checkpoints, attacks, defenses, a cross-model transferability matrix, and the interpretability tools across various robot embodiments and datasets.



\section{Broader Impacts}
This project will improve the security, robustness, and trustworthiness of foundation models in robotics, facilitating their safe and responsible deployment in critical systems. By addressing fundamental vulnerabilities and establishing comprehensive security frameworks, this work will increase public confidence and acceptance in robotic systems while cultivating a new generation of security-conscious researchers and practitioners in robotics and AI.
The project will provide open-source tools, datasets, benchmarks, and security frameworks, fostering collaboration and innovation in the research community and industry. These transformative resources will enable researchers and practitioners to evaluate and improve the security of robotic systems and accelerate progress in the field. The methodologies developed here will extend beyond robotics and benefit adjacent domains, contributing to the security of AI and multimodal models.
% Educational and outreach initiatives in this work will promote awareness of robotic and AI security challenges across academic and professional communities through curricula, labs, workshops, panel discussions, and seminars.
% This will create substantial knowledge-transfer pathways between academia and industry, ensuring that research advances are rapidly incorporated into professional practice.

\end{document}